
Epochs should be rejected regardless of blink presence when they contain raw-level unusable signal, such as electrode detachment, loose contact, motion bursts affecting the full segment, amplifier saturation, long flatlines, missing samples, or recording interruptions, because these conditions indicate that the data are fundamentally unreliable rather than merely influenced by normal ocular physiology. In addition, epochs should be dropped when the spatial pattern of the transient is inconsistent with true blink anatomy, for example when a large transient appears equally across all channels, when posterior channels dominate unexpectedly, or when the overall topography suggests movement burst, cable swing, or head-jerk artifact rather than a genuine blink. Another strong rejection criterion is flatness in the non-blink window: even if a disconnected sensor shows a blink-like spike due to noise pickup, the surrounding signal may reveal zero variance, near-zero peak-to-peak amplitude, or repeated constant values, all of which strongly suggest unusable data. High-frequency contamination is also important, because a true blink is usually smooth and low-frequency, so epochs should be rejected if the surrounding signal shows excessive high-frequency power, sharp repeated spikes, EMG-like activity, or unusually strong first-difference variance; this is especially useful in driving-task datasets for distinguishing real blinks from muscle bursts or steering-related movement. Finally, epochs should be excluded when they contain step artifacts or clipping, such as sudden baseline shifts, square-wave jumps, ADC saturation, repeated maximum or minimum values, or truncated waveform peaks, since a real blink should have a smooth onset and recovery rather than an abrupt non-physiological shape.
